%!TEX root = ../thesis.tex

\section{LLM-generated chapter version B: minimal additional simulations}
\label{sec:cpd-llm-version-b}

\noindent\textit{Draft note.}
This is an LLM-generated compact chapter draft for review. It keeps the Bernoulli BOCPD model as the formal core, but defines a small set of additional simulations that would make the story stronger. These simulations have not necessarily been run yet; the text states the expected qualitative outcome so that the required figures are defined in advance.

\subsection{Introduction}

Memories allow us to predict the future based on the past. This requires integrating previous experiences into an internal model of the world. Systems consolidation is usually discussed in this context as a way of extracting stable structure from individual experiences \citep{McClelland1995,Roxin2013}. However, statistical structure is not always stable. In a changing environment, old memories can become misleading, and the learner must decide whether past observations still belong to the same latent state as the present.

Bayesian change-point detection provides a normative formulation of this problem \citep{Adams2007}. The learner maintains uncertainty about the current run length, \(r_t\), the number of observations since the last change point. Short run lengths support rapid updating from recent observations; long run lengths support stable predictions based on many observations. The final prediction is a mixture over these possible evidence windows.

This chapter asks whether a pathway hierarchy from systems consolidation could provide a biological basis for such mixtures. PPT proposes that parallel pathways learn with different effective learning rates. This means that different pathways naturally estimate environmental statistics over different timescales. Here, we use Bayesian change-point detection to argue that this diversity may be useful for learning in changing environments. The stronger version of the chapter would not only show exact BOCPD, but also compare it to simpler learners and to finite pathway-like timescale banks.

\subsection{A Bernoulli change-point task}

We use a minimal binary task. On each trial,
\begin{equation}
    x_t \in \{0,1\}, \qquad x_t \sim \operatorname{Bernoulli}(\theta_t),
\end{equation}
where \(\theta_t\) is piecewise constant. At change points, \(\theta_t\) is redrawn from a prior distribution. This task is deliberately simple. It isolates the problem of estimating the current contingency without adding response categories, cues, or participant-level noise.

The stationary Bayesian learner assumes that \(\theta\) never changes. With prior \(\operatorname{Beta}(a_0,b_0)\), its posterior is
\begin{equation}
    \theta \mid x_{1:t} \sim \operatorname{Beta}\left(a_0+\sum_{i=1}^{t}x_i,\; b_0+\sum_{i=1}^{t}(1-x_i)\right).
\end{equation}
This learner becomes confident when the environment is stable, but adapts poorly after a change because old evidence remains in the posterior.

A fixed-learning-rate learner solves this only partly. It updates according to
\begin{equation}
    \hat\theta_t = \hat\theta_{t-1} + \alpha(x_t-\hat\theta_{t-1}).
\end{equation}
A large \(\alpha\) adapts quickly but is noisy during stable periods; a small \(\alpha\) is stable but slow after change points. This motivates the need for multiple timescales or an adaptive readout.

\begin{figure}[htbp]
\centering
\fbox{\begin{minipage}[c][0.32\textheight][c]{0.92\linewidth}
\textbf{Placeholder figure B1: why one timescale is not enough.}
Panel A: latent \(\theta_t\) and observations \(x_t\). Panel B: stationary Beta learner, showing slow recovery after a change. Panel C: fixed learning-rate learner with small \(\alpha\), stable but slow. Panel D: fixed learning-rate learner with large \(\alpha\), fast but noisy. Panel E: prediction error or negative log likelihood across learners.
\end{minipage}}
\caption{A changing Bernoulli contingency creates a stability--plasticity trade-off. A single evidence timescale cannot both average over stable periods and adapt rapidly after change points.}
\label{fig:cpd-timescale-tradeoff-placeholder-b}
\end{figure}

\subsection{Bayesian online change-point detection}

In Bayesian online change-point detection, the learner maintains a posterior over run lengths,
\begin{equation}
    p(r_t \mid x_{1:t}).
\end{equation}
The hazard function \(H(r)\) gives the probability that a run ends after surviving to age \(r\). The current implementation uses a constant hazard,
\begin{equation}
    H(r)=h,
\end{equation}
which corresponds to geometrically distributed regime durations.

For each possible run length, the Bernoulli parameter has a Beta posterior,
\begin{equation}
    \theta \mid r_t=r, x_{1:t}
    \sim \operatorname{Beta}\left(a_t^{(r)},b_t^{(r)}\right).
\end{equation}
The sufficient statistics \(a_t^{(r)}\) and \(b_t^{(r)}\) are the success and failure counts relevant under that run-length hypothesis, including prior pseudo-counts. The prediction under a run length is
\begin{equation}
    p(x_{t+1}=1 \mid r_t=r,x_{1:t})
    =\frac{a_t^{(r)}}{a_t^{(r)}+b_t^{(r)}}.
\end{equation}
The full prediction marginalizes over run lengths,
\begin{equation}
    p(x_{t+1}=1 \mid x_{1:t})
    =\sum_r \frac{a_t^{(r)}}{a_t^{(r)}+b_t^{(r)}}p(r_t=r\mid x_{1:t}).
\end{equation}

This model solves the timescale problem by not committing to one timescale. After evidence for a change point, posterior mass shifts toward short run lengths. During stable periods, posterior mass can move toward longer run lengths, allowing more evidence to be integrated.

\begin{figure}[htbp]
\centering
\fbox{\begin{minipage}[c][0.32\textheight][c]{0.92\linewidth}
\textbf{Placeholder figure B2: exact BOCPD simulation.}
Panel A: latent \(\theta_t\), observations, and true change points. Panel B: heat map of \(p(r_t\mid x_{1:t})\). Panel C: posterior mass near \(r_t=0\), interpreted as change-point probability. Panel D: BOCPD predictive mean. Panel E: comparison of predictive loss against stationary and fixed-learning-rate learners.
\end{minipage}}
\caption{Exact Bayesian online change-point detection tracks uncertainty over evidence windows. The expected result is that BOCPD adapts quickly after likely change points while preserving stable estimates during long runs.}
\label{fig:cpd-bocpd-placeholder-b}
\end{figure}

\subsection{Hazard assumptions}

The hazard function expresses assumptions about environmental stability. A constant hazard assumes that the next change is equally likely regardless of how long the current state has lasted. This is a good MVP because it is simple and standard. However, it is not the only possible assumption.

A decreasing hazard would express the belief that states which have already lasted for a long time are likely to continue. An increasing hazard would express the belief that regimes have typical lifetimes and become more likely to end as they age. A learned hazard treats volatility itself as uncertain \citep{Wilson2010}. These variants should be kept conceptually separate. A known age-dependent hazard is different from a learner that infers an unknown hazard.

For the minimal additional simulation, we do not need a full hazard-learning model. The useful comparison would be simpler: simulate the same Bernoulli sequence under constant hazard and evaluate observers with different assumed hazards. The expected result is that hazard mismatch changes the effective timescale of learning. High assumed hazard produces rapid but noisy updating; low assumed hazard produces stable but sluggish updating. This would make the link to subjective hazard and volatility concrete without adding a large model.

\begin{figure}[htbp]
\centering
\fbox{\begin{minipage}[c][0.30\textheight][c]{0.92\linewidth}
\textbf{Placeholder figure B3: hazard mismatch.}
Panel A: same generated sequence fit by BOCPD observers with low, correct, and high assumed hazard. Panel B: predictive means. Panel C: run-length posterior width or expected run length. Panel D: predictive loss as a function of assumed hazard. Optional Panel E: schematic of alternative hazard functions \(H(r)\): constant, decreasing, increasing.
\end{minipage}}
\caption{Assumed hazard controls the learner's effective timescale. The expected result is a trade-off between rapid adaptation and noise sensitivity.}
\label{fig:cpd-hazard-placeholder-b}
\end{figure}

\subsection{Finite timescale approximations}

Exact BOCPD maintains a posterior over many possible run lengths. This is statistically useful but computationally expensive. Several approximations reduce this problem to learning-rate control. Nassar et al. approximate the Bayesian observer with a single delta rule whose learning rate changes from trial to trial \citep{Nassar2010}. Wilson et al. instead approximate the Bayesian solution with a mixture of fixed-learning-rate estimators \citep{Wilson2013}.

The Wilson approximation is closest to the PPT story. Suppose we keep a finite set of estimates,
\begin{equation}
    \hat\theta_t^{(k)}, \qquad k=1,\ldots,K,
\end{equation}
with different fixed learning rates \(\alpha_k\), or equivalently different timescales \(\tau_k\). Each estimate is updated by
\begin{equation}
    \hat\theta_t^{(k)} = \hat\theta_{t-1}^{(k)} + \alpha_k(x_t-\hat\theta_{t-1}^{(k)}).
\end{equation}
A readout then combines them,
\begin{equation}
    \hat\theta_t = \sum_k w_t^{(k)}\hat\theta_t^{(k)}.
\end{equation}
The weights \(w_t^{(k)}\) play a role analogous to the run-length posterior. They should favor fast estimates after likely changes and slower estimates during stable periods.

A minimal simulation would compare three approximations: one fixed learning rate, a small hand-picked bank of fixed learning rates, and exact BOCPD. The expected result is that a finite bank improves strongly over one fixed learning rate and captures the main qualitative behavior of BOCPD, even if it does not match exact inference perfectly.

\begin{figure}[htbp]
\centering
\fbox{\begin{minipage}[c][0.34\textheight][c]{0.92\linewidth}
\textbf{Placeholder figure B4: finite pathway-like timescale bank.}
Panel A: estimates \(\hat\theta_t^{(k)}\) for fast, intermediate, and slow learners. Panel B: readout weights \(w_t^{(k)}\), shifting toward fast learners after change and slow learners during stability. Panel C: mixture prediction compared with exact BOCPD. Panel D: predictive loss for one fixed learning rate, finite bank, and exact BOCPD. Panel E: schematic mapping from \(\alpha_k\) or \(\tau_k\) to PPT pathway hierarchy.
\end{minipage}}
\caption{A finite bank of fixed-timescale learners can approximate the adaptive evidence-window selection performed by BOCPD. This is the simulation most directly connected to PPT.}
\label{fig:cpd-finite-bank-placeholder-b}
\end{figure}

\subsection{PPT as an approximation basis}

PPT proposes that different hippocampal--cortical pathways learn at different effective rates. In the original consolidation story, this produces different levels of semantic generalization. The change-point story adds another possible function: different learning rates produce estimates over different effective evidence windows.

This does not mean that each pathway literally represents one BOCPD run length. The safer claim is that pathways provide a basis. A fast pathway estimates recent contingencies. A slow pathway estimates stable regularities. Intermediate pathways fill in the range between them. A separate readout or control system can weight these estimates depending on uncertainty, surprise, inferred change points, or task demands.

This division of labor is important. Pathways supply candidate estimates; they do not by themselves compute the full run-length posterior. Neural systems implicated in volatility, uncertainty, and learning-rate control could provide part of the readout machinery \citep{Behrens2007,McGuire2014,Meder2017}. Thus, the biological proposal is distributed: pathway-specific learning rates provide the timescale basis, while frontal, cingulate, striatal, or neuromodulatory systems may control which timescale is trusted.

\begin{figure}[htbp]
\centering
\fbox{\begin{minipage}[c][0.30\textheight][c]{0.92\linewidth}
\textbf{Placeholder figure B5: conceptual model.}
Panel A: exact BOCPD trellis over run lengths \(r_t\). Panel B: compressed fixed-timescale basis with pathways indexed by \(k\). Panel C: readout/control signals selecting or weighting pathways. Panel D: thesis-level mapping: systems consolidation creates multiple semantic/timescale estimates; change-point inference provides a normative reason to keep and reweight them.
\end{minipage}}
\caption{Proposed mapping between Bayesian change-point inference and PPT. The claim is an approximation-basis claim, not a literal one-to-one mapping between pathways and run-length nodes.}
\label{fig:cpd-ppt-placeholder-b}
\end{figure}

\subsection{Discussion}

The central contribution of this chapter is the link between two ideas that are usually discussed separately. Bayesian change-point detection describes how an ideal learner should decide which past observations are still relevant. Systems consolidation describes how memories are transformed across hippocampal and cortical circuits. PPT suggests that these transformations occur across parallel pathways with different effective learning rates. The proposal here is that this pathway diversity may provide the fixed-timescale basis needed to approximate adaptive learning in changing environments.

The minimal simulations would serve different purposes. The exact Bernoulli BOCPD simulation demonstrates the normative target. The fixed-learning-rate comparison shows why one timescale is insufficient. The hazard-mismatch simulation shows that assumptions about volatility control the stability--plasticity trade-off. The finite-bank simulation is the key bridge to PPT: it would show that a small set of pathway-like estimators can reproduce the main qualitative benefit of the full run-length mixture.

There are clear limitations. First, the Bernoulli task is intentionally simple. The same conjugate structure generalizes to categorical observations with Dirichlet posteriors, but that is not needed for the MVP. Second, finite timescale banks require a readout mechanism. The pathway hierarchy provides estimates, not posterior weights by itself. Third, the model ignores replay. In biological systems, replay may change the effective evidence stream seen by each pathway, especially after surprising or boundary-related events. Finally, exact Bayesian optimality should not be overclaimed. The thesis claim is that PPT provides a plausible approximation basis for adaptive inference, not that it implements exact BOCPD.

\subsection{Chapter summary}

Changing environments require flexible control over evidence windows. BOCPD solves this by maintaining a posterior over run lengths and run-length-specific sufficient statistics \(a_t^{(r)}\) and \(b_t^{(r)}\). A finite mixture of fixed-timescale estimates can approximate this computation. PPT already proposes such a diversity of timescales in the form of parallel pathways with different effective learning rates. This provides a new normative reason for learning-rate diversity in systems consolidation.
